% !TeX root = closed-systems-from-comparison-completeness.tex
% ============================================================
\chapter{Relational Loop Defect and Riemann Curvature}
\label{ch:christoffel}
% ============================================================
\begin{chapterorientation}
\orientationitem{Input}{The discrete loop-defect package and the stitched refinement/observer arena \(ST\).}
\orientationitem{Role}{This chapter relates the algebraic commutator layer of transport to smooth curvature in a realized branch.}
\orientationitem{Output}{Riemann curvature as the smooth realization of relational loop defect, and the degree--\(2\) carrier used by the interface chapter.}
\end{chapterorientation}
\section{Introduction}
Under the standing principle of closed-world admissibility
(\cref{sp:closed-world}), this chapter develops the smooth-realization
consequences of the transport-visibility clause \cref{sp:transport}.
Combining the discrete obstruction arc of
\cref{ch:triangles,ch:descent_obstruction,ch:bridge} with the stitched
refinement/observer interleaving arena of \cref{ch:stitching}, it
derives the first smooth curvature realization used in
\cref{ch:interface}.
\Cref{ch:triangles,ch:descent_obstruction,ch:bridge} isolated an
intrinsic obstruction to
endpoint-determined transport.  At finite arity the first witness of
that obstruction occurs on triangle boundaries.  At the algebraic
level, its first visible carrier lies in the quadratic commutator
layer.  The purpose of the present chapter is to identify the canonical smooth
realization of that obstruction.
Differential geometry is not introduced here as an additional axiom.
The relevant obstruction has already been isolated on intrinsic
grounds: comparison structure, quotient semantics, transport, loop
defect, and triangle obstruction all arise before any smooth manifold
is mentioned.  The question is how this intrinsic obstruction appears
under smooth realization.
The answer is that for Levi--Civita transport on a smooth
pseudo-Riemannian manifold, the leading infinitesimal loop defect
around a minimal \(2\)-cell is precisely Riemann curvature; the
connection and curvature conventions are the standard differential
geometric ones \cite{KobayashiNomizu1963,Wald1984}.
Four ingredients enter.
First, the transport group carries the augmentation filtration, and the
first nontrivial commutator carrier is
\[
	F^2K/F^3K.
\]
Second, triangle boundaries furnish the first finite witnesses of
morphism-level transport defect.  The first obstruction is therefore
already intrinsically \(2\)-skeletal.
Third, triangle coherence is encoded by a quotient of the free path
groupoid characterized by a universal property: it is the initial
groupoid in which all directed triangle loops are trivial.
Fourth, under Levi--Civita realization, the leading infinitesimal
defect around a small \(2\)-cell is the curvature operator
\[
	R(\partial_i,\partial_j).
\]
The theorem proved in this chapter is that these structures match
canonically: the first intrinsic transport obstruction is quadratic and
\(2\)-skeletal, and its smooth infinitesimal realization is Riemann
curvature.
\begin{remark}[Ontological transition]
	Everything prior to smooth realization is intrinsic and algebraic:
	comparison structure, quotient semantics, transport, loop defect, and
	triangle obstruction.  Smooth manifolds enter only as realizations of
	the already isolated transport residue.  Geometry is therefore not
	postulated independently.  It is recognized as the smooth form of an
	intrinsic obstruction.
\end{remark}
% ------------------------------------------------------------
\section{Augmentation filtration and the first commutator carrier}
% ------------------------------------------------------------
Let \(K\) be a group.  Let
\[
	\varepsilon:\mathbb Z[K]\to \mathbb Z
\]
be the augmentation homomorphism, let
\[
	I:=\ker(\varepsilon)
\]
be the augmentation ideal, and let
\[
	\widehat{\mathbb Z[K]}
\]
denote the \(I\)-adic completion.
\begin{definition}[Augmentation filtration]
	For each \(m\ge 1\), define
	\[
		F^mK
		:=
		\{g\in K: g-1\in I^m\subseteq \widehat{\mathbb Z[K]}\}.
	\]
	The associated graded pieces are
	\[
		\gr^m K:=F^mK/F^{m+1}K.
	\]
\end{definition}
\begin{remark}[Degree \(1\) is tautological]
	\label{rem:F1-trivial}
	For every \(g\in K\), one has \(\varepsilon(g)=1\), hence
	\[
		g-1\in I.
	\]
	Therefore
	\[
		F^1K=K.
	\]
	Thus degree \(1\) carries no commutator information.  The first
	nontrivial graded transport residue can appear only in degree \(2\).
\end{remark}
\begin{lemma}[Commutators raise augmentation order]
	\label{lem:comm-raises}
	For all integers \(r,s\ge 1\),
	\[
		[F^rK,F^sK]\subseteq F^{r+s}K.
	\]
\end{lemma}
\begin{proof}
	Let \(a\in F^rK\) and \(b\in F^sK\).  Then there exist
	\[
		u\in I^r,\qquad v\in I^s
	\]
	such that
	\[
		a=1+u,\qquad b=1+v
	\]
	inside \(\widehat{\mathbb Z[K]}\).
	Since \(u\in I^r\) and \(v\in I^s\), the geometric-series expansions in
	the \(I\)-adic completion give
	\[
		a^{-1}=1-u+O(I^{r+1}),
		\qquad
		b^{-1}=1-v+O(I^{s+1}).
	\]
	Hence modulo \(I^{r+s+1}\),
	\begin{align*}
		aba^{-1}b^{-1}
		 & \equiv
		(1+u)(1+v)(1-u)(1-v) \\
		 & \equiv
		1+uv-vu.
	\end{align*}
	Since \(uv-vu\in I^{r+s}\), it follows that
	\[
		[a,b]-1\in I^{r+s}.
	\]
	Therefore
	\[
		[a,b]\in F^{r+s}K.
	\]
\end{proof}
\begin{corollary}[First visible noncommutativity]
	\label{cor:first-visible-comm}
	In a one-object transport system, the first graded carrier of transport
	noncommutativity is
	\[
		F^2K/F^3K.
	\]
\end{corollary}
\begin{proof}
	By \cref{rem:F1-trivial}, one has \(F^1K=K\).  Hence every commutator
	lies in \(F^2K\) by \cref{lem:comm-raises} with \(r=s=1\).  Quotienting
	by \(F^3K\) therefore isolates the first graded commutator residue.
\end{proof}
% ------------------------------------------------------------
\section{Triangle coherence and the universal triangle quotient}
% ------------------------------------------------------------
\Cref{thm:triangle,cor:bridge-triangle-regime} show that the first
finite witness of
morphism-level obstruction occurs on triangle boundaries.  We now
encode that fact categorically.
Let
\[
	\mathcal N=(V,E)
\]
be a finite directed graph, and let
\[
	\Path^\pm(\mathcal N)
\]
denote the free path groupoid on \(\mathcal N\).  Its objects are the
vertices \(V\), and its morphisms are reduced words in directed edges
and formal inverses.
\begin{definition}[Directed triangle loop (intrinsic form)]
	A \emph{directed triangle loop} in \(\Path^\pm(\mathcal N)\) is a loop
	of the form
	\[
		(p\to q)(q\to r)(r\to p)
	\]
	arising from a directed triangle in \(\mathcal N\). This is the same notion
	as \cref{def:bridge-triangle-loop}, written intrinsically in
	\(\Path^\pm(\mathcal N)\).
\end{definition}
\begin{definition}[Triangle congruence]
	Let \(\sim_\triangle\) denote the smallest groupoid congruence on
	\(\Path^\pm(\mathcal N)\) such that every directed triangle loop is
	identified with the identity at its basepoint.
\end{definition}
\begin{definition}[Triangle quotient groupoid]
	Define
	\[
		\Path^\pm_\triangle(\mathcal N)
		:=
		\Path^\pm(\mathcal N)/{\sim_\triangle},
	\]
	and write
	\[
		Q_\triangle:\Path^\pm(\mathcal N)\to \Path^\pm_\triangle(\mathcal N)
	\]
	for the quotient functor.
\end{definition}
\begin{theorem}[Universal property of the triangle quotient]
	\label{thm:triangle-quotient-universal}
	Let \(\mathcal G\) be any groupoid and let
	\[
		F:\Path^\pm(\mathcal N)\to\mathcal G
	\]
	be a functor such that for every directed triangle loop \(\ell\),
	\[
		F(\ell)=\id_{F(p)},
	\]
	where \(p\) is the basepoint of \(\ell\).  Then there exists a unique
	functor
	\[
		\overline F:\Path^\pm_\triangle(\mathcal N)\to \mathcal G
	\]
	such that
	\[
		F=\overline F\circ Q_\triangle.
	\]
	Equivalently, \(\Path^\pm_\triangle(\mathcal N)\) is initial among
	groupoids receiving a functor from \(\Path^\pm(\mathcal N)\) in which
	all directed triangle loops are trivial.
\end{theorem}
\begin{proof}
	By definition, \(\sim_\triangle\) is the smallest groupoid congruence
	generated by the requirement that every directed triangle loop be an
	identity, and closed under composition and inversion.  Since \(F\)
	sends each directed triangle loop to an identity and preserves
	composition and inversion, it is constant on
	\(\sim_\triangle\)-equivalence classes.  Therefore it factors through
	the quotient:
	\[
		F=\overline F\circ Q_\triangle
	\]
	for a unique functor
	\[
		\overline F:\Path^\pm_\triangle(\mathcal N)\to\mathcal G.
	\]
	Uniqueness follows because \(Q_\triangle\) is surjective on objects and
	morphisms.
\end{proof}
\begin{remark}[Structural meaning]
	The triangle quotient is not merely a convenient presentation.  It is
	the initial transport domain in which triangular \(2\)-cell boundaries
	are trivialized.  In that sense it is the canonical receptacle for the
	first \(2\)-skeletal coherence law.
\end{remark}
% ------------------------------------------------------------
\section{Single-object transport and based triangle defect}
% ------------------------------------------------------------
We now pass to the single-object transport regime.  Let
\[
	K
\]
be a group, viewed as a one-object groupoid.  A transport scheme on
\(\mathcal N\) is then an assignment of an element of \(K\) to each
directed edge, and by the universal property of the free path groupoid
it extends uniquely to a functor
\[
	\Hol:\Path^\pm(\mathcal N)\to BK,
\]
where \(BK\) is the one-object groupoid with automorphism group \(K\).
Fix a basepoint \(p\in V\).  Write
\[
	\Omega_p(\mathcal N):=\Hom_{\Path^\pm(\mathcal N)}(p,p)
\]
for the based loop group of the free path groupoid.
\begin{definition}[Based triangle normal subgroup]
	Let
	\[
		N_{\triangle,p}\trianglelefteq \Omega_p(\mathcal N)
	\]
	be the normal subgroup generated by all based conjugates of directed
	triangle loops, that is, by all loops of the form
	\[
		\alpha\,\ell\,\alpha^{-1},
	\]
	where \(\ell\) is a directed triangle loop based at some vertex \(q\)
	and \(\alpha:p\to q\) is any path in \(\Path^\pm(\mathcal N)\).
\end{definition}
\begin{definition}[Based triangle quotient loop group]
	Define
	\[
		\Omega_p^\triangle(\mathcal N)
		:=
		\Omega_p(\mathcal N)/N_{\triangle,p}.
	\]
\end{definition}
\begin{lemma}[Based loop group of the triangle quotient]
	\label{lem:triangle-based-loops}
	There is a canonical isomorphism
	\[
		\Hom_{\Path^\pm_\triangle(\mathcal N)}(p,p)
		\cong
		\Omega_p^\triangle(\mathcal N).
	\]
\end{lemma}
\begin{proof}
	Passing from \(\Path^\pm(\mathcal N)\) to the quotient
	\(\Path^\pm_\triangle(\mathcal N)\) imposes exactly the congruence
	generated by trivializing all directed triangle loops.  On the based
	loop group at \(p\), this identifies precisely the normal subgroup
	generated by all based conjugates of such loops.  Therefore the based
	loop group of the quotient is the quotient of \(\Omega_p(\mathcal N)\)
	by \(N_{\triangle,p}\).
\end{proof}
\begin{definition}[Based transport defect map]
	The transport functor \(\Hol\) induces a group homomorphism
	\[
		\delta_p:\Omega_p(\mathcal N)\to K
	\]
	by restriction to based loops at \(p\).
\end{definition}
\begin{theorem}[Triangle--loop identification]
	\label{thm:triangle-loop}
	The following are equivalent.
	\begin{enumerate}[label=\textup{(\roman*)}]
		\item
		      The transport functor \(\Hol\) factors through the triangle quotient
		      \[
			      Q_\triangle:\Path^\pm(\mathcal N)\to \Path^\pm_\triangle(\mathcal N).
		      \]
		\item
		      Every directed triangle loop is sent to the identity by \(\Hol\).
		\item
		      The subgroup \(N_{\triangle,p}\) lies in the kernel of
		      \[
			      \delta_p:\Omega_p(\mathcal N)\to K.
		      \]
	\end{enumerate}
	Equivalently, \(\Hol\) fails to factor through the triangle quotient if
	and only if the induced map on based loops descends to a nontrivial
	homomorphism
	\[
		\Omega_p^\triangle(\mathcal N)\to K.
	\]
\end{theorem}
\begin{proof}
	The equivalence of \textup{(i)} and \textup{(ii)} is exactly
	\cref{thm:triangle-quotient-universal} specialized to the functor
	\(\Hol\).
	For \textup{(ii)} \(\Longleftrightarrow\) \textup{(iii)}: if every
	directed triangle loop is sent to the identity, then every based
	conjugate of such a loop is also sent to the identity, since \(\delta_p\)
	is a homomorphism.  Hence the normal subgroup \(N_{\triangle,p}\) lies in
	\(\ker(\delta_p)\).  Conversely, every directed triangle loop based at
	\(p\) belongs to \(N_{\triangle,p}\), and every directed triangle loop at
	another basepoint becomes such after conjugation by a path from \(p\),
	so \(\ker(\delta_p)\supseteq N_{\triangle,p}\) implies that every
	directed triangle loop is sent to the identity.
	The final statement follows from \cref{lem:triangle-based-loops}: the
	failure of factorization is exactly the failure of \(\delta_p\) to kill
	\(N_{\triangle,p}\), equivalently the existence of a nontrivial induced
	map on the quotient
	\[
		\Omega_p^\triangle(\mathcal N)
		=
		\Omega_p(\mathcal N)/N_{\triangle,p}.
	\]
\end{proof}
\begin{remark}[Two-dimensionality of the first obstruction]
	\label{rem:triangle-2cell}
	Triangle loops are the first intrinsic \(2\)-cell boundaries in the
	transport system.  Their nontriviality detects the first finite failure
	of endpoint-determined transport.  The first obstruction is therefore
	already \(2\)-skeletal.
\end{remark}
% ------------------------------------------------------------
\section{Triangle boundaries and the quadratic commutator layer}
% ------------------------------------------------------------
We now relate the first \(2\)-skeletal obstruction to the quadratic
commutator carrier identified above.
\begin{lemma}[Triangle boundaries have trivial abelianized transport content]
	\label{lem:triangle-abelianization}
	Let
	\[
		\ab:\Omega_p(\mathcal N)\to \Omega_p(\mathcal N)^{\ab}
	\]
	be the abelianization map.  Then every element of \(N_{\triangle,p}\) has
	trivial image in the abelianized transport content relevant to
	degree \(1\).
\end{lemma}
\begin{proof}
	The subgroup \(N_{\triangle,p}\) is generated by based conjugates of
	directed triangle loops.  Passing to the abelianization kills
	conjugation and records only additive endpoint-level transport data.
	But a directed triangle loop is the boundary of a \(2\)-cell, and so
	carries no independent \(1\)-dimensional transport content.  Hence each
	generator of \(N_{\triangle,p}\) has trivial image in the abelianized
	transport content, and therefore so does every element of
	\(N_{\triangle,p}\).
\end{proof}
\begin{corollary}[Triangle boundary defect is invisible in degree \(1\)]
	\label{cor:triangle-degree1-invisible}
	Every defect supported on the subgroup
	\[
		N_{\triangle,p}\trianglelefteq \Omega_p(\mathcal N)
	\]
	has trivial image in the degree--\(1\) transport quotient
	\[
		F^1K/F^2K.
	\]
\end{corollary}
\begin{proof}
	By \cref{lem:triangle-abelianization}, every element of \(N_{\triangle,p}\)
	is invisible after passage to abelianized transport content.  By
	\cref{lem:gr1-endpoint}, the quotient
	\[
		F^1K/F^2K
	\]
	records exactly that abelianized endpoint transport content.  Hence any
	defect supported on \(N_{\triangle,p}\) has trivial image in
	\[
		F^1K/F^2K.
	\]
\end{proof}
\begin{lemma}[Degree--\(1\) detects only endpoint transport data]
	\label{lem:gr1-endpoint}
	The quotient
	\[
		F^1K/F^2K
	\]
	records only abelianized transport content of \(K\).  In particular, it
	cannot detect a defect supported on boundaries of \(2\)-cells.
\end{lemma}
\begin{proof}
	Since \(F^1K=K\) by \cref{rem:F1-trivial}, one has
	\[
		F^1K/F^2K=K/F^2K.
	\]
	By \cref{lem:comm-raises} with \(r=s=1\), every commutator in \(K\)
	lies in \(F^2K\).  Hence the quotient \(K/F^2K\) kills commutators and
	therefore factors through the abelianization \(K^{\ab}\).  It follows
	that \(F^1K/F^2K\) records only abelianized endpoint-level transport
	data and cannot detect boundary-supported \(2\)-cell defect.
\end{proof}
\begin{lemma}[Skeletal degree bound]
	\label{lem:2-skeletal-degree2}
	Let a transport defect be detected by loops supported on boundaries of
	\(2\)-cells.  Then its first nontrivial image in the augmentation
	filtration lies in \(F^2K\).
\end{lemma}
\begin{proof}
	By \cref{cor:triangle-degree1-invisible}, every defect supported on
	\(2\)-cell boundaries has trivial image in
	\[
		F^1K/F^2K.
	\]
	Therefore the first filtration level at which such a defect can appear
	is at least \(F^2K\).
\end{proof}
\begin{theorem}[First visible carrier of triangle obstruction]
	\label{thm:triangle-degree2}
	In a closed one-object transport regime, the first visible graded
	carrier of triangle obstruction is the quadratic quotient
	\[
		F^2K/F^3K.
	\]
\end{theorem}
\begin{proof}
	By \cref{thm:triangle-loop,rem:triangle-2cell}, triangle obstruction is
	the first finite witness of non-endpoint-determined transport and is
	already \(2\)-skeletal.  By \cref{lem:2-skeletal-degree2}, any such
	defect has trivial image in
	\[
		F^1K/F^2K
	\]
	and first possible nontrivial image in \(F^2K\).  By
	\cref{cor:first-visible-comm}, the first graded carrier of transport
	noncommutativity is precisely
	\[
		F^2K/F^3K.
	\]
	Hence the first visible graded carrier of triangle obstruction is
	\[
		F^2K/F^3K.
	\]
\end{proof}
\begin{remark}[Decategorified shadow]
	The quotient
	\[
		F^2K/F^3K
	\]
	is the first decategorified shadow of the first nontrivial
	\(2\)-dimensional transport obstruction.  The triangle quotient
	provides the \(2\)-skeletal source; the augmentation filtration
	provides the first graded algebraic image.
\end{remark}
% ------------------------------------------------------------
\section{Smooth realization by Levi--Civita transport}
% ------------------------------------------------------------
We now pass to smooth realization.  Let \((M,g)\) be a smooth
pseudo-Riemannian manifold with Levi--Civita connection \(\nabla\).
\begin{lemma}[Local geodesically convex neighborhoods]
	\label{lem:local-convex}
	For every \(x\in M\) there exists an open neighborhood \(U_x\) such
	that any two sufficiently near points of \(U_x\) are joined by a
	unique geodesic lying entirely in \(U_x\).
\end{lemma}
\begin{proof}
	Choose normal coordinates at \(x\).  For a sufficiently small
	star-shaped neighborhood of \(0\in T_xM\), the exponential map is a
	diffeomorphism onto its image.  In these coordinates geodesics are the
	images of radial straight lines in \(T_xM\).  Shrinking the image if
	necessary yields the required neighborhood.
\end{proof}
\begin{lemma}[Finite convex witness cover]
	\label{lem:finite-convex-cover}
	Let \(C\subseteq M\) be compact.  Then \(C\) admits a finite open cover
	\[
		\mathcal U=\{U_a\}_{a\in A}
	\]
	such that each \(U_a\) is geodesically convex and every nonempty finite
	intersection is geodesically convex after refinement.
\end{lemma}
\begin{proof}
	Choose neighborhoods \(U_x\) as in \cref{lem:local-convex} for each
	\(x\in C\).  Compactness gives a finite subcover.  Since only finitely
	many intersections occur, the sets may be shrunk simultaneously so
	that every nonempty finite intersection lies in a convex normal
	neighborhood and is therefore geodesically convex.
\end{proof}
\begin{definition}[Directed witness network]
	Fix a compact region \(C\subseteq M\) and choose a finite cover
	\[
		\mathcal U=\{U_a\}_{a\in A}
	\]
	as in \cref{lem:finite-convex-cover}.  Choose witness points
	\[
		p_a\in U_a
		\qquad(a\in A),
	\]
	and define the directed nerve network
	\[
		\mathcal N(\mathcal U)=(V,E)
	\]
	by
	\[
		V:=\{p_a:a\in A\},
		\qquad
		(a\to b)\in E \iff U_a\cap U_b\neq\varnothing,
	\]
	with orientation determined by a fixed total order on \(A\).
\end{definition}
\begin{definition}[Levi--Civita edge transport]
	\label{def:LC-transport}
	For each directed edge \(a\to b\) in the witness network, let
	\[
		\tau_{ab}:T_{p_a}M\to T_{p_b}M
	\]
	be parallel transport along the unique geodesic in
	\[
		U_a\cap U_b
	\]
	joining \(p_a\) to \(p_b\).
	These edge transports extend uniquely to a transport functor
	\[
		\Hol_{\LC}:\Path^\pm(\mathcal N(\mathcal U))\to \mathcal C^{\LC},
	\]
	where \(\mathcal C^{\LC}\) is the restriction of the frame groupoid of
	\(M\) to the witness vertices.
\end{definition}
% ------------------------------------------------------------
\section{Gauge invariance of the smooth loop defect}
% ------------------------------------------------------------
Fix a basepoint \(p\in V\).  The smooth loop defect is the homomorphism
\[
	\delta_p^{\LC}:\Omega_p(\mathcal N(\mathcal U))
	\to \Aut(T_pM).
\]
\begin{lemma}[Gauge conjugation of loop defect]
	\label{lem:conj}
	Let
	\[
		h_q\in \GL(T_qM)
		\qquad(q\in V)
	\]
	be a change of local frames.  Then the transformed loop defect at \(p\)
	satisfies
	\[
		\delta_p^{\LC\,\prime}(\gamma)
		=
		h_p\,\delta_p^{\LC}(\gamma)\,h_p^{-1}
		\qquad
		\text{for all }\gamma\in \Omega_p(\mathcal N(\mathcal U)).
	\]
\end{lemma}
\begin{proof}
	Under a frame change, each edge transport becomes
	\[
		\tau_{ab}'=h_b\,\tau_{ab}\,h_a^{-1}.
	\]
	For a based loop
	\[
		\gamma=e_n\cdots e_1
	\]
	at \(p\), the transformed holonomy is the product of these transformed
	edge maps.  The intermediate frame changes cancel telescopically,
	leaving only the basepoint terms:
	\[
		\delta_p^{\LC\,\prime}(\gamma)
		=
		h_p\,\delta_p^{\LC}(\gamma)\,h_p^{-1}.
	\]
\end{proof}
\begin{corollary}[Intrinsic smooth defect class]
	The conjugacy class
	\[
		[\delta_p^{\LC}]
	\]
	is independent of the choice of local frames.  Hence smooth loop defect
	determines an intrinsic conjugacy-class-valued transport invariant.
\end{corollary}
\begin{proof}
	Immediate from \cref{lem:conj}.
\end{proof}
% ------------------------------------------------------------
\section{\texorpdfstring{Curvature as infinitesimal \(2\)-cell defect}{Curvature as infinitesimal 2-cell defect}}
% ------------------------------------------------------------
We now compute the infinitesimal form of the smooth loop defect.
\begin{proposition}[Curvature commutator identity]
	Let \(\nabla\) be the Levi--Civita connection.  Then for every vector
	field \(V\),
	\[
		(\nabla_i\nabla_j-\nabla_j\nabla_i)V
		=
		R(\partial_i,\partial_j)V.
	\]
\end{proposition}
\begin{proof}
	Since the Levi--Civita connection is torsion-free, the coordinate
	vector fields satisfy
	\[
		[\partial_i,\partial_j]=0.
	\]
	By definition of curvature,
	\[
		R(\partial_i,\partial_j)V
		=
		\nabla_i\nabla_jV-\nabla_j\nabla_iV-\nabla_{[\partial_i,\partial_j]}V,
	\]
	and the final term vanishes.
\end{proof}
\begin{corollary}[Infinitesimal holonomy expansion]
	\label{cor:infdef}
	Let \(\Sigma_{ij}(\varepsilon)\) be an infinitesimal coordinate square
	of side length \(\varepsilon\) based at \(p\).  Then
	\[
		\Hol(\partial\Sigma_{ij}(\varepsilon))
		=
		\id
		+
		\varepsilon^2\,R(\partial_i,\partial_j)\big|_p
		+
		O(\varepsilon^3).
	\]
\end{corollary}
\begin{proof}
	Parallel transport along a coordinate edge in the \(i\)-direction has
	the expansion
	\[
		P_i=\id-\varepsilon \Gamma_i+O(\varepsilon^2),
	\]
	where \(\Gamma_i\) denotes the Christoffel matrix in the \(i\)-direction.
	Multiplying the four edge transports around the oriented square, the
	linear terms cancel.  The coefficient of \(\varepsilon^2\) is
	\[
		\partial_i\Gamma_j-\partial_j\Gamma_i
		+\Gamma_i\Gamma_j-\Gamma_j\Gamma_i,
	\]
	which is exactly the matrix of
	\[
		R(\partial_i,\partial_j).
	\]
	This gives the stated expansion.
\end{proof}
\begin{lemma}[Filtration degree and holonomy order]
	\label{lem:filtration-holonomy-order}
	Under smooth realization, the degree--\(m\) augmentation filtration
	corresponds to the order--\(m\) term in the small-loop holonomy
	expansion.  In particular, modulo \(F^3K\), only the quadratic
	infinitesimal holonomy term survives.
\end{lemma}
\begin{proof}
	Let \(\gamma_\varepsilon\) be a family of sufficiently small loops
	based at a fixed point, with size parameter \(\varepsilon\to 0\).  In a
	local trivialization, the corresponding holonomy admits an asymptotic
	expansion
	\[
		\Hol(\gamma_\varepsilon)
		=
		\id + A_1(\varepsilon)+A_2(\varepsilon)+A_3(\varepsilon)+\cdots,
	\]
	where \(A_m(\varepsilon)\) is homogeneous of order \(m\) in
	\(\varepsilon\).
	The augmentation filtration records the first nontrivial order at which
	deviation from the identity appears: membership in \(F^mK\) means that
	the holonomy differs from the identity only from order \(m\) onward.
	Hence passage to the graded quotient
	\[
		F^mK/F^{m+1}K
	\]
	isolates the order--\(m\) term in the small-loop expansion.
	For \(m=2\), modulo \(F^3K\) only the quadratic infinitesimal holonomy
	term survives.  Thus
	\[
		F^2K/F^3K
	\]
	identifies with the leading quadratic holonomy operator.
\end{proof}
\begin{corollary}[Quadratic holonomy survives modulo \(F^3K\)]
	\label{cor:quadratic-holonomy-modF3}
	Under smooth realization, the image of a sufficiently small loop in
	\[
		F^2K/F^3K
	\]
	is represented exactly by its quadratic holonomy term.
\end{corollary}
\begin{proof}
	This is the specialization of
	\cref{lem:filtration-holonomy-order} to \(m=2\).  Modulo \(F^3K\), all
	terms of order \(3\) and higher are discarded, and the quadratic term
	is the first surviving contribution.
\end{proof}
\begin{proposition}[Degree--\(2\) graded defect and infinitesimal holonomy]
	\label{prop:graded-holonomy}
	Under smooth realization, the degree--\(2\) graded piece
	\[
		F^2K/F^3K
	\]
	identifies with the leading infinitesimal holonomy operators around
	small \(2\)-cells.
\end{proposition}
\begin{proof}
	By \cref{cor:quadratic-holonomy-modF3}, the quotient
	\[
		F^2K/F^3K
	\]
	is represented under smooth realization by the quadratic term in the
	small-loop holonomy expansion.  For a small \(2\)-cell, that quadratic
	term is exactly the leading infinitesimal holonomy operator around the
	cell.
\end{proof}
\begin{proposition}[Christoffel integrability coefficients]
	\label{prop:christoffel}
	Let \(g\) be a smooth metric with Levi--Civita connection.  Then
	Christoffel's integrability coefficients coincide with the fully
	lowered Riemann curvature tensor:
	\[
		(gkhi)=R_{gkhi}.
	\]
\end{proposition}
\begin{proof}
	Christoffel's integrability coefficients are precisely the coordinate
	curvature coefficients of the Levi--Civita connection with indices
	lowered by the metric.  Substituting the Levi--Civita coefficients into
	the coordinate curvature formula and lowering the remaining index gives
	the identity.
\end{proof}
% ------------------------------------------------------------
\section{Relational loop defect realizes as curvature}
% ------------------------------------------------------------
We now combine the intrinsic algebraic carrier, the universal triangle
quotient, and the smooth transport computation.
\begin{theorem}[Relational loop defect realizes as curvature]
	\label{thm:christoffel-main}
	Under the standing principle \cref{sp:closed-world}, assume the
	Levi--Civita realization of \cref{def:LC-transport}. Then the
	infinitesimal relational loop defect is precisely Riemann curvature.
	More explicitly, for infinitesimal
	\(2\)-cells the leading nontrivial term of loop defect is
	\[
		R(\partial_i,\partial_j),
	\]
	and Christoffel's integrability coefficients are the corresponding
	lowered curvature components.
\end{theorem}
\begin{proof}
	Intrinsic transport theory identifies morphism-level obstruction with
	nontrivial based loop defect.  By
	\cref{thm:triangle-quotient-universal}, the triangle quotient is the
	initial transport domain in which all directed triangle loops vanish.
	Hence the first nontrivial finite obstruction is canonically
	\(2\)-skeletal.
	By \cref{thm:triangle-degree2}, the first visible graded carrier of
	that obstruction is the quadratic quotient
	\[
		F^2K/F^3K.
	\]
	By \cref{prop:graded-holonomy}, under smooth realization this graded
	carrier identifies with the leading infinitesimal holonomy operators
	around small \(2\)-cells.
	In the Levi--Civita realization, the latter are computed by
	\cref{cor:infdef}, whose leading nontrivial term is
	\[
		R(\partial_i,\partial_j).
	\]
	Therefore the smooth realization of the intrinsic infinitesimal
	\(2\)-cell defect is exactly Riemann curvature.
	The fully lowered form is Christoffel's integrability tensor, which
	coincides with the lowered Riemann tensor by
	\cref{prop:christoffel}.
\end{proof}
\begin{remark}[Jet compression versus holonomy detection]
	Curvature can be extracted in two distinct but compatible ways.
	First, holonomy detects curvature directly as transport defect around a
	small \(2\)-cell; this is the route intrinsic to the transport theory.
	Second, normal-coordinate jet expansion compresses the same curvature
	information into the quadratic coefficient of the metric.  The familiar
	factor \(\tfrac13\) belongs to the second route, not the first: it
	arises from symmetrization of the metric \(2\)-jet, not from the basic
	holonomy defect itself.
\end{remark}
\begin{remark}[Structural summary]
	This chapter identifies one transport-obstruction package through four
	linked descriptions:
	\[
		\begin{split}
			\text{triangle obstruction}
				 & \Longrightarrow \text{first \(2\)-skeletal transport defect}       \\
				 & \Longrightarrow \text{degree--\(2\) graded commutator carrier}     \\
				 & \Longrightarrow \text{smooth infinitesimal realization as Riemann curvature}.
		\end{split}
	\]
	Equivalently, Riemann curvature is the smooth infinitesimal
	realization of the first intrinsic transport noncommutativity, not a
	separate geometric input added from outside the closed transport
	chain.
\end{remark}
\begin{remark}[Position in the stack]
	This chapter identifies curvature as the smooth realization of the
	first intrinsic transport obstruction.  It does not yet pass to the
	stabilized quadratic carrier, the interface results
	(\cref{thm:interface-main,thm:interface-seal}), or the Einstein
	boundary analysis (\cref{sec:einstein-boundary}).  Those belong to the
	subsequent chapters.
\end{remark}

\section{Conclusion}
\label{sec:christoffel-conclusion}
The chapter's conclusion is exact: the first intrinsic noncommutative
transport defect, carried by \(F^2/F^3\), has Riemann curvature as its
smooth infinitesimal realization, with Christoffel's integrability
coefficients as its lowered coordinate form.

Accordingly, \cref{ch:interface} proves that the intrinsic quadratic
carrier is the first stable obstruction on the Stone limit and that,
under faithful smooth realization on that forced carrier, a nonzero
stabilized square class is equivalent to nonzero realized curvature.

\begin{chapterclosing}
\closingitem{Established}{The first intrinsic transport noncommutativity \(F^2/F^3\) realizes infinitesimally as Riemann curvature.}
\closingitem{Carried forward}{The next chapter stabilizes this degree--\(2\) carrier on the Stone-limit side and analyzes faithfulness.}
\end{chapterclosing}
